% Appendix: the P-series record, 7 to 23 Oct 2026. Draft 27 Nov 2026 (writer id: pseries).
% Every number and dated fact is traced in app_pseries_NOTES.md. Summarised in Chapter ch:legacy and Table tab:T21; this appendix
% gives the per-experiment record and does not change any value or verdict given there.
\chapter{The P-series: diagnostics and semi-supervised learning (7 to 23 Oct)}\label{app:pseries}

This appendix records, experiment by experiment, the work of 7 to 23 Oct 2026: the diagnostics P1 to P7 that rebuilt the supervised
baseline, the semi-supervised (SSL) series P8 to P17 with the audit that interrupted it, the exploratory P18, the per-cell
evaluation E1 and its unseen-module extension, and the full ALHVR adaptation AL1. Chapter~\ref{ch:legacy} gives the short account
and the lessons of the period, Table~\ref{tab:T21} the verdict rows and Appendix~\ref{app:t21} the arm codes. Nothing here changes a
value or a verdict given there; this appendix adds the protocol, the arms, the result and the consequence of each experiment.
Tables~\ref{tab:pseries-sum1} and~\ref{tab:pseries-sum2} at the end list all of them.

\section{Conditions of every number in this appendix}\label{sec:pseries-conditions}

\paragraph{Data and labels.} All runs use the legacy split of the BenchmarkEL \texttt{20211104} release (478 train, 54 validation, 50
test cells, split seed 0). P1 to P5 examined models or probes trained on the 20211104 masks (v0 of Section~\ref{sec:data-labels}); from P6 on every run trains and
scores on the revised E-SCDD masks \citep{hanifi2025escdd}, label set v1, except where P14 trains on v0 on purpose. A label
fraction of 5, 10, 20, 50 or 100\% means 23, 47, 95, 239 or 478 labelled train cells; inactive area occurs in 2, 4, 9 and 20 of the
first four sets. The in-domain unlabelled pool is the rest of the train split (455, 431, 383, 239 cells); the external pool holds
11,938 CSIR and ARTS cells, 77\% CSIR and 23\% ARTS.

\paragraph{Metrics.} Unless stated, IoU is the per-class IoU of crack or of inactive, pooled over all pixels of the 50 test cells
(one intersection and one union per class over the set), at the argmax prediction of the final weights; it is not a class mean,
and background does not enter it. relF1@2 is the pooled relaxed F1 at 2~px of Section~\ref{sec:data-f1}. AP is pixel average
precision from score histograms with 1,000 bins. P6 and P7 decided on crack AP and inactive AP; from P8 on the three primaries are
crack IoU, inactive IoU and crack relF1@2. All values are on a 0 to 1 scale.

\paragraph{Decision rule.} Arms were compared by the locked rule of Chapter~\ref{ch:legacy} (P8 rule 1): a difference is BETTER or
WORSE if it exceeds $\max(0.01,\, 2\sqrt{(s_A^2 + s_B^2)/3})$, and a comparison passes if one primary is better and none is worse.
At 3 against 3 seeds this is about 2 standard errors, roughly a one-sided $p \approx 0.06$; the divisor stays 3 at 5 against 5 seeds,
where it is about 2.58 standard errors. In tables a value is written as the difference followed by the rule's threshold in parentheses.
Three robustness checks were reported beside each verdict and decided nothing: a paired image bootstrap over the 50 test cells
(2,000 resamples, 95\% interval), the difference without the 3 cells with the most labelled pixels of the class
(``top-3 out''), and a split by source (21 CSIR and ARTS cells, 29 CFVS, BMRK and SDLE cells).

\paragraph{Training.} The reference recipe from P6 on (L3) is U-Net++ \citep{zhou2018unetpp} with an ImageNet ResNet-34 encoder
\citep{he2015resnet,iakubovskii2019smp}, cross-entropy plus soft Dice on crack and inactive, label-safe augmentation (flips, quarter
turns, brightness and contrast), 20,000 steps, labelled batch 8, Adam at $10^{-3}$ with weight decay $10^{-4}$, 200 warm-up steps and
a cosine decay, final weights. SSL arms add an unlabelled batch of 8 and a Mean Teacher \citep{tarvainen2017meanteacher} with
EMA 0.999, whose teacher is scored; the unlabelled loss weight is 0 for 2,000 steps and rises linearly to 1 over the next 2,000.
Most runs used the training machine (RTX A6000). P1 to P5, the supervised references of P10, the replication and calibration arms
of P12 and the six fine-tuning runs of P15 used the second machine (P15's pretraining run used the training machine). Seeds of the
two machines were never pooled, and no verdict rule compared arms across machines; the P12 calibration contrast compared the two
machines' supervised arms on purpose and decided nothing.

\paragraph{What these numbers can and cannot show.} The limits of Section~\ref{sec:legacy-conditions} apply to every value below.
test50 was scored in every run, overlaps the training pool (Section~\ref{sec:data-split}) and has 3 cells that hold 59\% of its
inactive pixels; every test value is in-distribution with known leakage, context only. No P-series model was selected on test
(final weights; validation-only choices), but the set was seen after every run. E1 values in this period were read without
registration. ``Passes'' means a pass of the lenient locked rule, not a confirmation under the protocol of Chapter~\ref{ch:protocol}.

\section{Restart and the diagnostics P1 to P5 (7 to 10 Oct)}\label{sec:pseries-diag}

On 7 Oct the work restarted with analysis before any new method. Three doubts drove it: a supervised model scored about 0.3 crack
IoU on the 47 images it had trained on, the old unlabelled loss seemed to have no effect, and the published numbers did not match
ours. Every protocol from here on was committed before its runs (the first commit at 23:59 on 7 Oct), and anything outside a
protocol was labelled exploratory.

\emph{Conditions of P1 to P5.} These five experiments examined models, labels or training images of the old pipeline: the from-scratch U-Net++ with
weighted cross-entropy and strong augmentation, the SemiCE Mean Teacher and the old ALHVR port, trained on the 20211104 masks on
the legacy split, whose test cells were scored in every run and which overlaps the training pool, as later found. Their numbers are
reproduction and diagnostic records of that pipeline, not valid results, and none of them is comparable with any later value.

\paragraph{P1: overfit check (8 Oct).} Question: why did a supervised model reach only 0.325 crack IoU on its own 47 training
images, while inactive reached 0.88? Ten crack-bearing training images were fitted for 3,000 steps by five arms that separate class
weights, augmentation, loss and numeric precision (one run each, three seeds for the two precision arms), scored on the same ten
images. Without augmentation, weighted cross-entropy reached 0.612 crack IoU (AP 0.929), plain cross-entropy 0.677, cross-entropy
plus soft Dice 0.861; with the training augmentation, weighted cross-entropy reached 0.381 (AP 0.587). Verdicts of the locked rules:
augmentation is the largest single cap (H2 met); class weights cost precision (0.62 against 0.81) but are not the main cap (H1 not met
as written, partial support); label inconsistency does not cap fitting (H3 refuted); no pipeline bug, since loader masks equal the
raw files in 10 of 10 images and fp32 against bf16 gave 0.656 $\pm$ 0.048 against 0.668 $\pm$ 0.031 over three seeds (H4 refuted).
It led to label-safe augmentation and to a Dice term on held-out data (steps L1 and L3 of P6).

\paragraph{P2: SemiCE gradient probe (8 Oct).} Question: why did the old Mean Teacher at $\lambda = 1.5$ score the same as its
$\lambda = 0$ control? A forward and backward probe on 40 training batches of a final old Mean Teacher checkpoint (no optimiser
step) measured the unlabelled gradient. The locked rule fired: 100\% of unlabelled pixels had at least two ``positive'' classes
(threshold 20\%). The gradient was neither wasted on a common logit shift (share 0.041 $\pm$ 0.009 against 0.5) nor negligible
(0.25 $\pm$ 0.35 of the supervised norm), but it pushed background, padding and crack up together. Every old SemiCE Mean Teacher
result is therefore a test of that loss, not of Mean Teacher, and is recorded as invalid. A side measurement: a standard
confidence rule (teacher probability $\geq$ 0.95) covered 77\% of unlabelled pixels, of which 0.024\% were crack, about 8 times
below crack's labelled share. This rare-class starvation became the central SSL difficulty; a correct hard pseudo-label loss
\citep{sohn2020fixmatch} replaced SemiCE, and the queued external-pool SemiCE runs were held and never started.

\paragraph{P3: label versions (10 Oct, exploratory).} The revised E-SCDD masks were found on 10 Oct. Re-scoring final old models on
both versions (no training, no protocol) moved test crack IoU by $+0.05$ to $+0.07$ for every model, for example 0.304 to 0.355
for the 10\% model (3 seeds). Of 281 crack images, 129 had a changed crack mask, 34 of them in test, and the changes lay almost
entirely within 2~px of the old labels: a cleaning (dotted labels joined and widened), not a new annotation. From P6 on every run
trained and scored on the revised masks.

\paragraph{P4: frozen-feature probe (10 Oct).} Linear probes on frozen features (trained on 47 cells with the old masks, 2,000 steps,
two probe seeds; scored on pooled validation and test cells with revised labels) gave crack AP 0.029 for a random ResNet-50, 0.321 for ImageNet ResNet-50, 0.331 for
ImageNet ConvNeXt-T and 0.183 for DINOv2 ViT-S/14 \citep{oquab2023dinov2}; inactive AP 0.097, 0.423, 0.383 and 0.384.
Locked rule 1 fired (best pretrained above random by $+0.30$ crack AP), so masked-image pretraining and foundation-teacher
distillation stayed on the list; rule 2 fired (DINOv2 below ImageNet ResNet-50 on crack), so ViT patch features were dropped for
crack. It led to the ImageNet encoder of step L2.

\paragraph{P5: pseudo-labels inside ALHVR's regions (10 Oct).} ALHVR \citep{lei2025alhvr} routes unlabelled pixels by the
confidence of two networks into regions where both, one or neither is sure. On two old ALHVR runs and 378 unlabelled train cells,
crack pseudo-labels in the ``neither sure'' region had strict precision 0.27 and relaxed (2~px) precision 0.64; at the confident end
(ensemble probability $\geq$ 0.95) 0.59 and 0.96, but that rule covered only 8\% of ground-truth crack. Under per-class thresholds the
relaxed precision in that region was 0.48 (recall 0.29), so locked rule 1 fired (borderline, below 0.5): hard crack pseudo-labels
are unusable there. Rule 2 fired for inactive (strict precision 0.43); per-class thresholds moved 40 points of crack out of the
region but cost precision elsewhere (rule 3 not met). Our reading: crack pseudo-labels find the right crack but miss its exact
pixels. This was the measured motivation for the tolerant pseudo-label loss: a rare-class pixel of the teacher is credited if the
student predicts that class anywhere within 2~px, confident head-class pixels away from rare targets get cross-entropy, and other
pixels get no loss.

\section{The corrected baseline P6 and the pilot P7 (10 Oct)}\label{sec:pseries-p6}

\paragraph{P6: a one-change ladder.} Locked at 11:42 on 10 Oct; 21 runs (seeds 0 to 2) on the training machine. Steps at 10\%
labels: L0 the old recipe (from-scratch 5M U-Net++, strong augmentation, weighted cross-entropy), L1 label-safe augmentation, L2 the
ImageNet ResNet-34 U-Net++, L3 cross-entropy plus Dice; L0 also at 100\%, L3 also at 20\% and 100\%. Table~\ref{tab:pseries-p6}
gives the results. Rule 1 on test AP: L0 to L1 helps both classes (crack AP $+0.020$, threshold 0.019; inactive AP $+0.028$,
threshold 0.020); L1 to L2 and L2 to L3 hurt both; L0 from 10\% to 100\% has no detectable effect; L3 from 10\% to 100\% helps.
Rule 2 picked the base of P7 by validation: L0 (score 0.480). Rule 3 measured L0's label headroom as $-0.001$ pooled AP, below
0.02, so P7 became a pilot. An interim headroom of $+0.026$ had averaged the per-split APs and was replaced by the pooled value.

Two exploratory findings mattered more than the verdicts. First, AP was not comparable across these changes: with the ImageNet
encoder and Dice many labelled pixels score below 0.001, where AP cannot rank them (labelled crack pixels below 0.001: 3\% for L0,
19\% for L2, 42\% for L3), while argmax crack IoU rose from L1 to L2 and stayed level to L3. From P8 the primaries became argmax
IoU and relF1@2. Second, only L3 turned labels into scores: from 10\% to 100\% its crack IoU rose by $+0.075$ $[+0.029, +0.128]$ and
inactive IoU by $+0.135$ $[+0.049, +0.205]$, while L0 gained nothing measurable. SSL can only close a gap that exists, so L3 became
the supervised reference ``A'' of every later SSL experiment. A later correction (12 Oct): most of the L1 minus L0 crack IoU gain
($+0.048$) came from one L0 seed that put 71,376 false crack pixels on the noisiest test cell; without that cell the gain is about
$+0.010$.

\begin{table}[htbp]
\centering\small
\caption{P6 ladder, test50, label set v1, mean $\pm$ sd over seeds 0 to 2 (training machine). Pooled per-class IoU and AP, argmax,
final weights. In-distribution, known leakage, context only; not comparable with the protocol tables.}
\label{tab:pseries-p6}
\begin{adjustbox}{max width=\linewidth}
\begin{tabular}{@{}llllll@{}}
\toprule
Arm & crack AP & inactive AP & crack IoU & inactive IoU & crack relF1@2 \\
\midrule
L0 10\% (old recipe) & 0.553 $\pm$ 0.015 & 0.621 $\pm$ 0.003 & 0.319 $\pm$ 0.051 & 0.439 $\pm$ 0.015 & 0.746 $\pm$ 0.078 \\
L1 10\% (label-safe aug.) & 0.574 $\pm$ 0.006 & 0.649 $\pm$ 0.017 & 0.368 $\pm$ 0.004 & 0.473 $\pm$ 0.014 & 0.806 $\pm$ 0.004 \\
L2 10\% (ImageNet ResNet-34) & 0.519 $\pm$ 0.003 & 0.477 $\pm$ 0.018 & 0.392 $\pm$ 0.003 & 0.383 $\pm$ 0.024 & 0.806 $\pm$ 0.003 \\
L3 10\% (CE + Dice) & 0.400 $\pm$ 0.004 & 0.302 $\pm$ 0.018 & 0.396 $\pm$ 0.003 & 0.294 $\pm$ 0.019 & 0.814 $\pm$ 0.003 \\
L3 20\% & 0.433 $\pm$ 0.004 & 0.335 $\pm$ 0.042 & 0.398 $\pm$ 0.002 & 0.347 $\pm$ 0.043 & 0.822 $\pm$ 0.008 \\
L0 100\% & 0.545 $\pm$ 0.003 & 0.612 $\pm$ 0.006 & 0.337 $\pm$ 0.004 & 0.391 $\pm$ 0.008 & 0.782 $\pm$ 0.008 \\
L3 100\% & 0.593 $\pm$ 0.003 & 0.420 $\pm$ 0.018 & 0.471 $\pm$ 0.001 & 0.429 $\pm$ 0.011 & 0.876 $\pm$ 0.003 \\
\bottomrule
\end{tabular}
\end{adjustbox}
\end{table}

\paragraph{P7: tolerant pseudo-labels on L0 (pilot).} At 10\% labels on the L0 base, three Mean Teacher arms ran with seeds 0 to 2:
B hard pseudo-labels (confidence $\geq$ 0.95), C tolerant pseudo-labels, D the C pipeline with the unlabelled weight at 0. A detail
that mattered later: the labelled images received the student's stronger photometric view. Verdicts (test AP): C minus D crack
$+0.000$ (n.d.) and inactive $-0.032$ (worse), so C does not help; B minus D worse on both; C minus B better on both ($+0.024$ and
$+0.064$). Width check: C's relaxed precision rose ($+0.031$), so its cracks were thinner, not fatter. All of it is a pilot, since L0
had no label headroom. The only contrast robust over cells was C over B on inactive IoU, $+0.056$ $[+0.022, +0.086]$. It led to P8.

\section{The SSL series on L3 with a shared batch-normalisation forward: P8 to P12 (10 to 11 Oct)}\label{sec:pseries-joint}

From P8 the base is L3 and the primaries are crack IoU, inactive IoU and crack relF1@2. The base choice was made after seeing P6
and is disclosed as post hoc; it applies equally to every arm.

\paragraph{P8: P7's arms on L3, 10\% labels (locked 15:43, 10 Oct).} No unlabelled loss passed against its loss-off control: C minus
D crack IoU $+0.009$ (threshold 0.010), inactive $+0.059$ (0.132), relF1@2 $+0.015$ (0.017); B minus D and C minus B were n.d. as
well. The discovery was the pipeline handicap: the loss-off control D was below plain supervised A, robustly over cells, by
$-0.013$ $[-0.028, -0.001]$ crack IoU and $-0.101$ $[-0.171, -0.031]$ inactive IoU. Scoring the student instead of the EMA teacher gave
the same numbers, so the teacher was not the cause. Three differences remained: the stronger labelled view, batch-normalisation
statistics computed over 8 labelled and 8 unlabelled images in one forward pass (``joint BN''), and the random data streams.

\paragraph{P9: the 11,938-cell external pool, 10\% labels (10 to 11 Oct).} Question: was P8's null a matter of pool size? With the
external pool, Cx minus Dx passed on inactive IoU ($+0.056$, threshold 0.046) with nothing worse, while Bx did not help and Cx did
not beat Bx. The pass was the same size as P8's C minus D ($+0.059$) and passed only because seed spread was smaller; over cells its
interval included 0, $[-0.013, +0.117]$, and it halved without the top 3 cells. The pool-size effect on the loss was $-0.005$,
$-0.003$ and $-0.007$ on the three primaries. Two exploratory findings: swapping the pool moved the loss-off control (Dx minus D
inactive $+0.094$ $[+0.003, +0.179]$), although with the loss off the pool can act only through batch statistics and data streams; and
the external pool lowered crack on test cells of its own sources (Cx minus C on CSIR and ARTS cells $-0.019$ $[-0.041, -0.003]$).

\paragraph{P10: 5\% labels, both pools (10 to 11 Oct).} No loss passed. In-domain, C5 and B5 (added by an amendment before any B5
run) were n.d. against D5, with seed standard deviations near 0.1: seed 2 collapsed under both pseudo-label losses (crack IoU 0.142
and 0.128) but not under its loss-off control, and in that seed the tolerant arm's rare-class pseudo-label share stayed at 0.001 to 0.002 from step
4,000, a collapse visible in the training log without labels. With the external pool, Cx5 minus Dx5 was worse on relF1@2
($-0.027$, threshold 0.024) and fattened cracks (relaxed precision $-0.065$); its crack IoU difference, $-0.012$ $[-0.026, -0.001]$, lay
on the pool's own sources ($-0.065$ $[-0.092, -0.026]$ on the 21 CSIR and ARTS cells). The supervised references ran on the second
machine and entered no rule; descriptively, the loss-off control sat far below them (D5 minus A5 crack IoU $-0.055$).

\paragraph{P11: matched labelled view, 10\% labels (locked 00:11, 11 Oct).} The labelled images now received exactly the supervised
arm's augmentation; batch normalisation stayed joint. The matched view removed most of P8's handicap (Dm minus D inactive $+0.101$),
leaving a small crack handicap, Dm minus A $-0.007$ $[-0.016, -0.001]$, n.d. by the seed rule. The loss helped: Cm minus Dm passed on
inactive IoU ($+0.071$, threshold 0.069) and relF1@2 ($+0.018$, threshold 0.013), robust over cells ($+0.071$ $[+0.020, +0.121]$ and
crack $+0.009$ $[+0.003, +0.016]$). SSL beat supervised: Cm minus A passed on inactive IoU ($+0.071$, threshold 0.067) and relF1@2
($+0.013$, threshold 0.012), and was credited because the loss also passed against its own control. Over cells only the relF1@2 part
was robust, $+0.013$ $[+0.005, +0.024]$; the inactive part was not, $+0.071$ $[-0.009, +0.158]$. Cmx minus A passed the arithmetic on
inactive but was not credited, because Cmx minus Dmx failed. P11 at 10\% labels is the one passed and credited ``SSL beats
supervised'' result under a shared batch-normalisation forward.

\paragraph{P12: matched view at 5\%, with a supervised reference on the same machine (locked 01:37, 11 Oct).} Cm5 minus Dm5 passed on
crack IoU ($+0.023$, threshold 0.022) and relF1@2 ($+0.040$, threshold 0.027) and replicated on the second machine with fresh seeds
3 to 5 ($+0.047$ and $+0.078$). Cm5 minus A5c failed: $+0.006$ (threshold 0.014), $+0.003$ (0.161) and $+0.015$ (0.022). The loss-off
control was below supervised on crack, $-0.017$ $[-0.024, -0.010]$, n.d. by the seed rule; the matched view had already removed the
labelled-view part (Dm5 minus D5 $+0.045$). Our reading at the time: the pseudo-label loss mostly offset a handicap of the joint
forward, consistent with batch-normalisation mixing. The second machine's arms passed the arithmetic against its own supervised
reference but, by the locked rule, could not rescue the failed rule on the training machine. A calibration contrast found the two
machines' supervised 10\% arms $+0.004$ apart on crack IoU (n.d.).

\section{The audit and P13 (11 to 13 Oct)}\label{sec:pseries-p13}

On 11 Oct (11:24) the supervised baseline still beat SSL in most contrasts, and an audit asked whether the implementation or the
design was at fault. The code checks found no bug: teacher and student views share geometry, the EMA covers the batch-normalisation
buffers, the teacher is detached and in evaluation mode, the teacher is scored, and no unlabelled image is in test. The audit found
two deviations from the methods that report gains. The student's perturbation was weak: the unlabelled loss stayed about 10 times
below the supervised loss (0.004 against 0.038 at step 20,000 in one 5\% run) while the teacher was confident on 94 to 97\% of
unlabelled pixels; segmentation consistency is reported to need mask-based mixing such as CutMix \citep{yun2019cutmix}
\citep{french2020semiseg}. And batch normalisation
was joint. Recomputing the batch-normalisation statistics of finished checkpoints on all 478 train cells lifted every arm by about
$+0.01$ crack IoU but left the D5 minus A5 gap at about $-0.06$: the handicap was in the trained weights, so the forward pass during
training had to change. In 100-step GPU smoke runs the student with a separate unlabelled forward (statistics frozen for unlabelled
images) and loss weight 0 was bit-identical to supervised training (0 of 350 tensors differing), so from P13 the supervised run is
the exact loss-off control.

\paragraph{P13 protocol.} Locked at 11:51 on 11 Oct before any training run; Amendment 1 at 12:51, before any P13 run, restricted it
to 5\% labels, dropped the loss-off arm (identical to supervised) and the second-machine arms, and gave the primary contrast five
seeds, by the standing rule of at least 5 seeds for effects below 0.05 and a one-day deadline. Arms: A5c supervised (seeds 0 to 4,
three from P12), Cs5 tolerant pseudo-labels with the separate forward (3 seeds), CsM5 Cs5 plus CutMix at probability 0.5 (5 seeds),
BsM5 hard pseudo-labels with the separate forward and CutMix (3 seeds). No value was tuned. 13 runs on the training machine, 11 Oct
13:30 to 21:04.

\paragraph{P13 results.} The primary CsM5 minus A5c passed on 5 against 5 seeds: inactive IoU $+0.117$ (threshold 0.052), relF1@2
$+0.033$ (0.025), crack IoU $+0.014$ (0.016, n.d.); both passing metrics were robust over cells, $[+0.042, +0.189]$ and
$[+0.006, +0.055]$, and survived top-3 removal ($+0.103$, $+0.020$). Secondary: BsM5 minus A5c passed on all three ($+0.036$, $+0.113$,
$+0.052$; crack $[+0.021, +0.048]$); Cs5 minus A5c passed on inactive only ($+0.105$); the CutMix effect CsM5 minus Cs5 was n.d. by the
seed rule ($+0.010$ crack IoU) though robust over cells, $[+0.004, +0.017]$; tolerant minus hard under CutMix, CsM5 minus BsM5, was
worse on crack IoU ($-0.022$, threshold 0.014); separate minus joint, Cs5 minus Cm5, was n.d. (Cm5's inactive sd was 0.139), with
inactive $+0.084$ $[+0.030, +0.143]$ over cells. Every SSL seed's test inactive IoU (0.318 to 0.403) lay above every supervised
seed's (0.192 to 0.284), while validation, with only 6 inactive cells, did not show the gain. No arm collapsed. CutMix raised the
unlabelled-to-supervised loss ratio at step 20,000 (seed 0) from 0.46 (Cs5) to 0.77 (CsM5) and 0.98 (BsM5).

\paragraph{Amendment 2 (13 Oct, post hoc).} Because BsM5's crack gain ($+0.036$) was below 0.05 on 3 seeds, BsM5 received seeds 3
and 4 with a stopping rule fixed in advance (5 seeds in total, reported whatever they show). Both came in lower (0.317 and 0.327
crack IoU). On 5 against 5, BsM5 minus A5c still passed on all three, $+0.029$ (0.016), $+0.107$ (0.044), $+0.045$ (0.023), crack
$[+0.014, +0.041]$ over cells, and CsM5 minus BsM5 turned from worse to n.d. ($-0.015$, threshold 0.018), still below 0 over cells.
Without the noisiest test cell (post hoc) CsM5 minus A5c crack IoU crosses its threshold ($+0.016$ against 0.014).

\section{P14 and P15: label noise and encoder transfer (11 Oct)}\label{sec:pseries-p14}

\paragraph{P14: offset-tolerant supervision.} Locked at 12:55 on 11 Oct; 9 runs at 10\% labels, seeds 0 to 2, training machine.
The tolerant rule was applied to human labels (rare-class labels credited within 2~px, the Dice term kept exact) on old (v0) and
revised (v1) masks, all scored on v1. Rule 1: the label-noise gap L3 minus L3o exists, $+0.082$ crack IoU (threshold 0.011) and
$+0.081$ relF1@2, robust over cells ($[+0.046, +0.129]$) and in both source groups. Rule 2: tolerance recovered none of it (L3To minus
L3o $-0.004$; at least $+0.041$ was needed): fails. Rule 3: tolerance on revised labels passed on inactive IoU ($+0.037$, threshold
0.028), not robust over cells ($[-0.016, +0.093]$), with crack unchanged and crack AP lower (0.400 to 0.380). Scored on the old
labels, the revised-label model still scored higher (crack IoU 0.351 against 0.332). Our reading: the old and revised labels differ by more
than an offset a tolerant loss can absorb; the revision changed what is labelled, not only where.

\paragraph{P15: UCF-EL encoder pretraining.} Locked at 12:55 on 11 Oct. A DeepLabV3+ \citep{chen2018deeplabv3plus} with an ImageNet
ResNet-34 encoder was trained on all 9,562 UCF-EL train cells \citep{fioresi2022ucfel} (59,760 steps, 100 epochs); its 216 encoder
tensors initialised our U-Net++ with a fresh decoder. No UCF-EL cell reached a near-duplicate correlation of 0.90 with our validation
or test cells (maximum 0.79). Six runs on the second machine, seeds 0 to 2, each against that machine's ImageNet arm: at 5\%, gU5
minus A5 was n.d. on all three (crack $+0.0096$ against 0.0100); at 10\%, gU10 minus gA10 was worse on all three, $-0.012$ (0.011),
$-0.086$ (0.059), $-0.014$ (0.013), robust over cells. Later qualification (30 Oct, Section~\ref{sec:legacy-pf}): under the house
fine-tune the deep encoder is erased in every ResNet-34 U-Net++ run measured, so in our reading P15 tested little more than the
transfer of the first encoder layers; this was not measured on the P15 encoder.

\section{P16: 20, 50 and 100\% labels (11 to 12 Oct)}\label{sec:pseries-p16}

We asked on 11 Oct why SSL had only been run at 5 and 10\%, where inactive occurs in 2 or 4 labelled cells and the teacher is
weakest. Part A (supervised references) was locked at 17:28, before any P16 run, and Part B (SSL arms) at 17:35, before any P16
SSL run. The recipe was fixed
by a rule on P13 \emph{validation} only: CutMix on if CsM5 scored above Cs5 on the mean of the three validation primaries. CsM5 scored
0.3507 and Cs5 0.3652, so CutMix was off and P16 ran the Cs recipe (tolerant pseudo-labels, separate forward, matched view), which on
P13 test was the weakest SSL arm on crack; test was not used for the choice. Arms (5 seeds each, training machine): S20 and S50 with
the in-domain pool (383 and 239 cells), Sx100 with the external pool (there is no in-domain pool at 100\%), each against its supervised
reference (seeds 0 to 2 of P6 pooled with new seeds 3 and 4 at 20\% and 100\%, after a pooling check).

All three primaries failed, n.d. on all three metrics: S20 minus A20 $-0.040$ (0.069), $+0.029$ (0.052), $-0.070$ (0.111), with a
width flag; S50 minus A50 $+0.039$ (0.088), $+0.006$ (0.034), $+0.053$ (0.135); Sx100 minus A100 $+0.028$ (0.091), $+0.003$ (0.025),
$+0.040$ (0.122). Every bootstrap interval included 0. The large thresholds came from one test cell: in 4 of the 24 runs (2 of 15
supervised, 2 of 15 SSL) the model labelled 15 to 31\% of the noisiest cell of all 582 as crack, 39k to 80k false pixels against a
true crack area of 787~px, which lowered the run's crack IoU by 0.11 to 0.18. A post hoc check without that cell, not a locked rule,
gave crack differences of $+0.002$, $+0.002$ and $-0.002$, all n.d.; the verdict does not change. At these fractions the unlabelled
loss was 2 to 7\% of the supervised loss at step 20,000, against 46\% for Cs5 at 5\%. Our reading: for the recipe that was run, the
inactive gain seen at 5\% shrank as labels grew and was not detectable from 20\% up.

\section{E1: per-cell metrics and the unseen-module extension (11 to 13 Oct)}\label{sec:pseries-e1}

Pooled IoU is dominated by a few cells and depends on the operating point, while the intended use needs per-cell quantities. E1
definitions were fixed at 12:59 on 11 Oct, before any E1 number: crack tolerant F1 at 2~px (mean over cells with at least 20 crack
pixels), crack length error from skeletons, crack false alarm (share of crack-free cells with at least 50 predicted crack pixels),
inactive area error, inactive detection F1 and inactive cell IoU (per cell, mean over cells with at least 0.5\% inactive area). Two
cell sets were scored: test50 and 106 cells of release \texttt{20221008} that are not among our 582 cells and not in the external pool,
from 67 modules of which only one has cells among our 582 (101 CSIR, 5 BMRK; 32 with crack, 10 with inactive). These 106 cells are
the E1 set of Section~\ref{sec:data-sets}. A disclosed amendment at 13:16, made after the first-pass extension numbers were seen,
set four defect types that our taxonomy lacks (scuff, corrosion, brightening, star) to ignore; the first pass had counted crack
predictions on them as false alarms.

Against the criterion ``adopt as primary if the per-cell interval is narrower'', per-cell crack metrics were not more sensitive than
pooled IoU or relF1@2 on test (contrasts excluding 0: 8 for tolerant F1 against 13 for crack IoU), so the pooled metrics stayed
primary; inactive cell IoU separated more contrasts than pooled inactive IoU (9 against 7 on test, 17 against 10 on the extension).
The extension is a different regime: L3 at 100\% fell from crack IoU 0.471 on test to 0.152, and from tolerant F1 0.80 to 0.41;
L3 at 10\% from 0.396 to 0.080. Table~\ref{tab:pseries-e1} lists the extension reads attached to the experiments; all were
secondary, decided nothing, and were read on the same 106 cells many times (Section~\ref{sec:legacy-taught}).

\begin{table}[htbp]
\centering\small
\caption{E1 extension reads of the P-series (106 unseen-module cells, per-cell metrics unless stated, argmax, final weights,
mean over seeds, paired cell bootstrap). Secondary, unregistered, repeated reads of the same cells; the E1 label set of
Section~\ref{sec:data-labels} (release 20221008 masks, four defect types set to ignore). Pooled IoU here is per-class IoU pooled over the 106 cells.}
\label{tab:pseries-e1}
\begin{adjustbox}{max width=\linewidth}
\begin{tabular}{@{}lll@{}}
\toprule
Experiment, contrast & Metric & Difference [95\%] \\
\midrule
P11 Cm $-$ A (10\%) & crack false alarm & $+0.113$ $[+0.049, +0.186]$ (more false alarms) \\
P11 Cm $-$ A & inactive cell IoU & $+0.071$ $[+0.001, +0.166]$ \\
P12 Cm5 $-$ A5c (5\%) & crack false alarm & $+0.095$ $[+0.035, +0.162]$ (more false alarms) \\
P12 Cm5 $-$ A5c & inactive cell IoU & $+0.212$ $[+0.065, +0.351]$ \\
P13 CsM5 $-$ A5c (5\%) & inactive cell IoU & $+0.467$ $[+0.322, +0.592]$ (0.033 to 0.500) \\
P13 CsM5 $-$ A5c & crack IoU, pooled & $+0.001$ $[-0.011, +0.012]$ \\
P13 BsM5 $-$ A5c & crack relF1@2, pooled & $+0.046$ $[+0.018, +0.072]$ \\
P13 Cs5 $-$ Cm5 (separate $-$ joint BN) & crack false alarm & $-0.095$ $[-0.141, -0.050]$ \\
P14 L3 $-$ L3o (revised $-$ old labels) & crack IoU, pooled & $+0.030$ $[+0.012, +0.048]$ \\
P14 L3T $-$ L3 & inactive cell IoU & $+0.088$ $[+0.005, +0.193]$ \\
P15 gU5 $-$ A5 & crack tolerant F1 & $-0.034$ $[-0.068, -0.004]$ \\
P15 gU10 $-$ gA10 & every metric & n.d. \\
P16 S20 $-$ A20 & inactive cell IoU & $+0.179$ $[+0.011, +0.385]$ \\
P16 S50 $-$ A50 & crack relF1@2, pooled & $+0.047$ $[+0.019, +0.077]$ \\
P16 Sx100 $-$ A100 & every metric & n.d. \\
P17 BsD5 $-$ BsM5 & crack false alarm & $+0.049$ $[+0.011, +0.089]$ \\
AL1 AL5 $-$ A5 & every metric & interval includes 0 \\
\bottomrule
\end{tabular}
\end{adjustbox}
\end{table}

Our reading of the extension reads: under the shared batch-normalisation forward, tolerant pseudo-labels raised crack false alarms
on unseen modules by about 0.10 of crack-free cells and improved inactive there; the separate forward removed most of the false-alarm
rise, and separate-forward SSL lifted inactive on unseen modules strongly (CsM5 $-$ A5c inactive cell IoU $+0.467$).

\section{P17, P18 and the checks of 13 Oct}\label{sec:pseries-p17}

\paragraph{P17: a decoupled decoder.} Question: does BsM5's gain need the labelled decoder to fit pseudo-labels, or does it enter
through the encoder? Of the two methods we considered on 13 Oct, one was a 3D method with three expert decoders and was not run; from
SKCDF \citep{zhang2025skcdf} only its data-flow decoupling was taken,
the component its own ablation credits with the largest share of its gain. In the arm BsD5 (BsM5 plus the switch, seeds 0 to 4, locked
at 01:51 on 13 Oct) the student's unlabelled forward runs through a copy of the decoder that follows the labelled decoder by an
exponential average of 0.99, so the unlabelled loss reaches only the encoder and that copy; pseudo-labels and scoring still use the
EMA teacher. A built-in check passed: up to step 2,000, where the unlabelled weight is 0, every logged line of each BsD5 seed equalled
BsM5's. Primary BsD5 minus BsM5: $-0.006$ (0.013), $+0.011$ (0.043), $-0.008$ (0.017), n.d.; the thresholds are 44\%, 41\% and 39\% of
the BsM5 minus A5c gain, so a decoder share below about 40\% of that gain cannot be detected with 5 seeds. BsD5 minus A5c passed on
all three ($+0.023$, $+0.117$, $+0.037$). Over cells, crack leaned negative, $-0.006$ $[-0.011, +0.002]$, with the lean in the CSIR
and ARTS cells. BsD5's teacher ranked pixels less well (crack AP 0.310 against 0.367). A follow-up check found the EMA teacher and
the student of BsM5, CsM5 and BsD5 within $\pm$0.005 on every metric at step 20,000, so the SSL gains are carried by the trained
network itself.

\paragraph{P18: the BsM recipe at 10 and 20\% (exploratory).} BsM was chosen after seeing P13's test cells, so P18 was labelled
exploratory, although its protocol was locked at 17:31 on 13 Oct before any run. BsM10 (3 seeds) minus A10 passed the arithmetic on
inactive only ($+0.132$, threshold 0.032; $[+0.043, +0.235]$ over cells); BsM20 (5 seeds) minus A20 on inactive only ($+0.060$,
threshold 0.051; $[-0.018, +0.119]$ over cells), with one seed failing on the noisiest cell. Our reading: P16's ``no gain from 20\% up''
depends at least partly on the recipe, at least for inactive; which of BsM's two differences from Cs matters was not identified. No
locked verdict changed. On 14 Oct a published multi-class EL segmentation model \citep{hanifi2025multiclass} was reproduced with its
released code (P19); that reproduction is a baseline record, not part of the SSL series.

\section{21 to 23 Oct: twin test cells, SK1 and AL1}\label{sec:pseries-al1}

From 21 Oct the remaining SSL questions were run by a small team of AI agents in which a reviewer agent, separate from the agent
running the arms, had to approve each protocol before launch and recomputed every reported number from the raw per-run files.

\paragraph{Twin test cells and a reduced pool.} The data review of 21 Oct found no newer labelled release and kept the legacy split
with v1. It fixed a list of 7 test cells that are the same physical cell as a train cell (same source and module position, high-pass
correlation 0.80 to 0.94, identical crack geometry); they hold 36,405 of 92,074 test crack pixels (39.5\%) and 12.1\% of test inactive
pixels. Section~\ref{sec:data-split} gives the later, fuller count of shared cells. At 5\% every train twin was in the
\emph{unlabelled} pool, so the 5\% SSL arms had trained on those cells and the supervised arm had not. Every later table was reported
with and without the 7 cells and the noisiest cell, both lists labelled post hoc relative to the P-series. On re-scoring P13 and P17
without the twins, BsM5 minus A5c still passed on all three, crack only just ($+0.019$, threshold 0.018); CsM5 minus A5c passed on
inactive only; BsD5 minus BsM5 stayed n.d. We decided the same day that new protocols drop the 7 train twins from the unlabelled pool
(455 to 448 cells) and re-run the SSL references, and that the earlier claims stay, with the ``without twins'' line beside them.

\paragraph{SK1 and a fixed logit bias (21 to 22 Oct).} SK1 added SKCDF's auxiliary balanced head to the decoupled recipe at 5\% on
the reduced pool (5 seeds per arm). The primary, BsDB5 minus BsD5, was n.d. in all four cell sets ($-0.0044$ (0.0100), $-0.0113$
(0.0282), $+0.0044$ (0.0100)), and the SKCDF track was closed on 22 Oct. A post hoc inference-only check then subtracted the log class
prior of each run's labelled cells from its logits ($\tau = 1$, fixed before any number): this alone raised crack IoU by $+0.015$ to
$+0.023$ at 5\%. With the bias given to both arms, the SSL crack gain survived (BsM5 $+0.037$); with the bias given to the supervised
arm only, BsM5's crack gain over it was n.d. ($+0.0135$, threshold 0.0159). This yardstick says the 5\% SSL crack gain is no larger
than what a free bias gives the supervised model; it is not a verdict against SSL, and the inactive gains survived it.

\paragraph{AL1: full ALHVR adapted to EL (protocol 21 Oct, read 23 Oct).} Question: at 5\% labels, does ALHVR \citep{lei2025alhvr},
adapted to 2D EL cells on the L3 base with a schedule under which its unlabelled loss is actually applied, beat step-matched
supervised training? The adaptation changed the old port in steps each grounded in a diagnostic run before the protocol: the L3 base,
the P-series loss schedule, a separate batch-normalisation forward, feature noise only (the reference feature dropout zeroed 96 to
99.7\% of encoder positions on this base), prototypes from the full-resolution decoder output (at the reference 1/8 resolution the two
networks' prototypes were orthogonal, which explains the old port's 4.6\% crack accuracy), and class-normalised sharpening with a
detached target. The maximum loss weight was matched to BsM5's unlabelled gradient share: 0.1 for ALHVR (the paper's value), 1.0 for
cross pseudo-supervision \citep{chen2021cps}. Arms, 5 seeds each on the training machine, all
SSL arms on the reduced pool: A5 supervised, AL5 ALHVR-EL, ALd5 without its prototype loss, ALdh5 that arm at weight 0.3, CPS5 plain
cross pseudo-supervision on all pixels, CPSm5 the same with ALHVR's ``both sure'' region masked out (the one-change routing test), and
BsM5 re-run as a reference. Checks: network A at weight 0 equalled supervised tensor for tensor; A5 equalled the earlier A5c seed for
seed; all 35 runs reached step 20,000 (launched 21 Oct 15:20, last rescore 23 Oct).

Table~\ref{tab:pseries-al1} gives the locked verdicts. The primary AL5 minus A5 was WORSE on all three primaries, with paired
cell-bootstrap intervals below 0 on all three, and WORSE in all four cell sets (all 50, without the noisiest cell, without the 7 twins,
without both). The routing test and the prototype-loss test were n.d.; every ALHVR-family arm was WORSE than supervised on inactive
IoU; cross pseudo-supervision was n.d. against supervised and below BsM5. In a post hoc reading, the ALHVR-family losses more than
halved inactive recall (A5 0.276, AL5 0.124), while cross pseudo-labels raised it (CPS5 0.400). Four weak predictions written before
any AL1 score was opened all held; the inactive damage was not predicted. Our reading: on this base ALHVR's losses are active and
hurt, neither of its modules adds anything in a one-change contrast, and plain two-network cross pseudo-supervision is level with
supervised. This closed the ALHVR question that had run through the project since before 7 Oct; from 24 Oct the work turned to models
trained on all labels (Chapter~\ref{ch:legacy}).

\begin{table}[htbp]
\centering\small
\caption{AL1 locked verdicts, 5\% labels, 5 against 5 seeds, test50 (all 50 cells), label set v1, training machine. Difference
(rule threshold) for crack IoU / inactive IoU / crack relF1@2; pooled per-class IoU, argmax, final weights, network A (EMA teacher for BsM5). SSL arms on the
448-cell pool. In-distribution, known leakage, context only.}
\label{tab:pseries-al1}
\begin{adjustbox}{max width=\linewidth}
\begin{tabular}{@{}llll@{}}
\toprule
Contrast & Role & Difference (threshold) & Verdict \\
\midrule
AL5 $-$ A5 & primary & $-0.0169$ (0.0131) / $-0.1324$ (0.0761) / $-0.0312$ (0.0216) & WORSE on all three \\
CPSm5 $-$ CPS5 & routing & $-0.0025$ / $-0.0125$ / $-0.0026$ & n.d. \\
AL5 $-$ ALd5 & prototype loss & $-0.0079$ / $+0.0169$ / $-0.0167$ & n.d. \\
ALd5 $-$ A5 & secondary & inactive $-0.1493$ (0.0631) & WORSE \\
ALdh5 $-$ A5 & secondary & inactive $-0.1056$ (0.0539) & WORSE \\
CPS5 $-$ A5 & secondary & $+0.0034$ / $+0.0736$ (0.0737) / $+0.0068$ & n.d. \\
AL5 $-$ CPSm5 & method & $-0.0178$ / $-0.1935$ / $-0.0354$ & WORSE \\
BsM5 $-$ A5 & descriptive & $+0.0274$ / $+0.1320$ / $+0.0412$ & passes \\
\bottomrule
\end{tabular}
\end{adjustbox}
\end{table}
Robustness of the primary (reported, deciding nothing): crack IoU $[-0.0287, -0.0065]$, inactive IoU $[-0.1906, -0.0750]$, relF1@2
$[-0.0536, -0.0118]$.

\section{Summary of the P-series}\label{sec:pseries-summary}

\begin{table}[htbp]
\centering\footnotesize
\caption{The P-series, part 1: P1 to P10 (7 to 11 Oct 2026). P1 to P5 examine old-pipeline models and labels (old pipeline, test set
scored in every run, overlapping splits as later found; history only). P6 to P10: legacy split, label set v1, test50 primaries
(in-distribution, known leakage, context only). Verdicts are those of each experiment's locked rule. TM, SM: training, second machine.}
\label{tab:pseries-sum1}
\begin{adjustbox}{max width=\linewidth}
\begin{tabular}{@{}>{\raggedright\arraybackslash}p{0.07\linewidth}>{\raggedright\arraybackslash}p{0.17\linewidth}>{\raggedright\arraybackslash}p{0.21\linewidth}>{\raggedright\arraybackslash}p{0.30\linewidth}>{\raggedright\arraybackslash}p{0.17\linewidth}@{}}
\toprule
Exp. & Date; labels; seeds; machine & Question & Result and verdict & Led to \\
\midrule
P1 & 8 Oct; 10 train images; 1 run per arm (3 for two); SM & why crack does not fit its own training images & no bug; augmentation is the main cap; CE + Dice fits to 0.861 & label-safe aug.\ (L1), Dice (L3) \\
P2 & 8 Oct; probe, 40 batches; SM & does SemiCE move the student & rule fired: 100\% of pixels with $\geq$ 2 positives; old SemiCE results invalid & hard pseudo-label loss \\
P3 & 10 Oct; re-score; exploratory; SM & does the label version matter & crack IoU $+0.05$ to $+0.07$ for every model & revised masks from P6 \\
P4 & 10 Oct; probe, 2 probe seeds; SM & do frozen features carry the signal & rules 1 and 2 fired; ImageNet CNN best; DINOv2 weaker on crack & ImageNet encoder (L2) \\
P5 & 10 Oct; offline; SM & are pseudo-labels usable in ALHVR's regions & rules 1, 2 fired (not usable in the unsure region); rule 3 not met & tolerant pseudo-labels \\
P6 & 10 Oct; 10, 20, 100\%; 3; TM & ladder L0 to L3 & L1 helps; L2, L3 hurt AP; only L3 has label headroom; base by rule = L0, headroom $-0.001$ & P7 a pilot; L3 the reference \\
P7 & 10 Oct; 10\%; 3; TM & tolerant pseudo-labels on L0 & pilot; C $-$ D n.d. / worse; C beats B & P8 \\
P8 & 10 Oct; 10\%; 3; TM & the same on L3 & no loss passes; loss-off control below supervised & P9 to P11 \\
P9 & 10 to 11 Oct; 10\%, external pool; 3; TM & does pool size matter & Cx $-$ Dx passes on inactive ($+0.056$), not robust over cells & P10, P11 \\
P10 & 10 to 11 Oct; 5\%; 3; TM (+SM) & does SSL help more at 5\% & no loss passes; seed collapse; external pool worse on relF1@2 & P12, P16 \\
\bottomrule
\end{tabular}
\end{adjustbox}
\end{table}

\begin{table}[htbp]
\centering\footnotesize
\caption{The P-series, part 2: P11 to AL1 (11 to 23 Oct 2026). Legacy split, label set v1, test50 primaries (in-distribution, known
leakage, context only); P18 is exploratory; SK1 and AL1 use the 448-cell pool. TM, SM: training, second machine.}
\label{tab:pseries-sum2}
\begin{adjustbox}{max width=\linewidth}
\begin{tabular}{@{}>{\raggedright\arraybackslash}p{0.07\linewidth}>{\raggedright\arraybackslash}p{0.17\linewidth}>{\raggedright\arraybackslash}p{0.21\linewidth}>{\raggedright\arraybackslash}p{0.30\linewidth}>{\raggedright\arraybackslash}p{0.17\linewidth}@{}}
\toprule
Exp. & Date; labels; seeds; machine & Question & Result and verdict & Led to \\
\midrule
P11 & 11 Oct; 10\%; 3; TM & matched labelled view & Cm $-$ Dm passes; Cm $-$ A passes and is credited (inactive $+0.071$, relF1@2 $+0.013$) & P12, P13 \\
P12 & 11 Oct; 5\%; 3; TM (+SM) & the same at 5\% & Cm5 $-$ Dm5 passes, replicates on SM; Cm5 $-$ A5c fails & the audit, P13 \\
P13 & 11 Oct (Am.\ 2 13 Oct); 5\%; 5 v 5; TM & separate BN forward and CutMix & primary passes: inactive $+0.117$, relF1@2 $+0.033$, crack IoU n.d.; BsM5 passes on all three & P16, P17 \\
P14 & 11 Oct; 10\%; 3; TM & tolerant loss on human labels & gap $+0.082$ exists; recovery fails; inactive pass on v1 not robust & label quality as the larger lever \\
P15 & 11 Oct; 5, 10\%; 3; SM & UCF-EL encoder pretraining & n.d. at 5\%; WORSE on all three at 10\% & closed for this fine-tune \\
P16 & 11 to 12 Oct; 20, 50, 100\%; 5 v 5; TM & SSL with more labels & n.d. at all three fractions & P18 \\
P17 & 13 Oct; 5\%; 5 v 5; TM & decoupled decoder & n.d.; a share below about 40\% not excluded & nothing launched \\
P18 & 13 Oct; 10, 20\%; 3, 5; TM & BsM at more labels (exploratory) & inactive passes the arithmetic at both & no locked change \\
SK1 & 21 to 22 Oct; 5\%; 5 v 5; TM & balanced head on the decoupled recipe & n.d. in all four cell sets & SKCDF closed \\
AL1 & 21 to 23 Oct; 5\%; 5 v 5; TM & full ALHVR-EL & WORSE on all three primaries in all four cell sets & ALHVR closed \\
\bottomrule
\end{tabular}
\end{adjustbox}
\end{table}

\section{What the record adds to Chapter~\ref{ch:legacy}}\label{sec:pseries-adds}

The lessons of the period are those of Section~\ref{sec:legacy-taught}. The per-experiment record adds three points of method. First,
a loss-off control can carry a handicap of its own: from P8 to P12 every unlabelled loss was compared with a control that sat below
plain supervision, so ``the loss helps'' and ``SSL beats supervised'' had to be separate rules, and only P11 passed both under the
shared forward. Second, in our reading the largest measured movements came from the labels and the pipeline rather than from the form of the
loss: training on revised rather than old labels moved crack IoU by $+0.082$ (P14), and the matched labelled view moved the
loss-off control's inactive IoU by $+0.101$ (P11), while tolerant and hard pseudo-labels were n.d. against each other at 5 against 5
seeds (P13, Amendment 2). Third, the 50 test cells set the limits of every verdict: three cells decide inactive, one noisy cell can
wreck a run, and the seven test cells of the 21 Oct list had trained-on twins in the 5\% unlabelled pool. These are the reasons the protocol of
Chapter~\ref{ch:protocol} develops on module-grouped cross-validation of the pool.
